\begin{titlepage}
\thispagestyle{empty}
\begin{tikzpicture}[remember picture,overlay]
\fill[Ink] (current page.south west) rectangle (current page.north east);
\begin{scope}[shift={(current page.south west)},x=1mm,y=1mm]
\clip (0,0) rectangle (180,255);
\draw[white,opacity=.055,step=8mm] (15,32) grid (188,173);
\foreach \t in {-4,-3,...,4}{
  \draw[Teal!45!white,opacity=.35,line width=.5pt,domain=-5:5,samples=180,smooth,variable=\x]
  plot ({100+16*\x},{104+1.25*(\x*\x*\x-5*\x)+\t*5});
}
\draw[Gold,line width=1.15pt,domain=-4.5:4.5,samples=170,smooth,variable=\x]
plot ({100+16*\x},{104+1.25*(\x*\x*\x-5*\x)});
\foreach \x/\y in {64/117,100/104,136/91}{\fill[Paper] (\x,\y) circle (1.0mm);\draw[Gold,opacity=.6] (\x,\y) circle (2.2mm);}
\draw[Gold,line width=.7pt] (23,215)--(42,215);
\node[anchor=west,text=Gold,font=\sffamily\fontsize{10}{12}\selectfont] at (46,215) {GROTHENDIECK\enspace /\enspace GAETA};
\node[anchor=north west,text=Paper,align=left,font=\sffamily\bfseries\fontsize{30}{39}\selectfont] at (22,202) {函子代数几何\\导论};
\node[anchor=west,text=Paper,font=\sffamily\fontsize{16}{20}\selectfont] at (23,163) {第一卷\quad 仿射代数几何};
\node[anchor=west,text=Paper!65!Ink,align=left,font=\fontsize{9}{14}\selectfont] at (23,147) {INTRODUCTION TO FUNCTORIAL\\ALGEBRAIC GEOMETRY\\[2pt]VOL. I\enspace AFFINE ALGEBRAIC GEOMETRY};
\node[text=Paper!85!Ink,font=\fontsize{14}{17}\selectfont] (a) at (42,57) {$\mathfrak X_A$};
\node[text=Paper!85!Ink,font=\fontsize{14}{17}\selectfont] (b) at (137,57) {$\Hom_{k\text{-}\mathrm{alg}}(A,-)$};
\draw[-{Stealth[length=2mm]},Paper!60!Ink,line width=.5pt] (a)--node[above,text=Gold,font=\small] {$\sim$} (b);
\draw[Gold,opacity=.6,line width=.4pt] (23,37)--(157,37);
\node[anchor=west,text=Paper!80!Ink,font=\sffamily\fontsize{9}{12}\selectfont] at (23,27) {A. Grothendieck 讲授\qquad F. Gaeta 笔记};
\node[anchor=west,text=Paper!55!Ink,font=\sffamily\fontsize{8}{10}\selectfont] at (23,19) {中文译本\quad /\quad 中·英·法术语互链};
\end{scope}
\end{tikzpicture}
\null
\end{titlepage}
